\begin{table*}[t]
\centering
\caption{Characteristics of the 17 BRGM piezometric series and four public multivariate forecasting datasets. BRGM statistics summarize the 17 separate series; all seasonality, trend and residual-noise scores are computed on the chosen training target.}
\label{tab:dataset_characteristics}
\footnotesize
\setlength{\tabcolsep}{4pt}
\begin{tabular}{@{}lcccccc@{}}
\toprule
Dataset & Sampling & $M$ & $T$ & Seasonality & Trend & Residual ratio \\
\midrule
BRGM piezometers ($n=17$) & 1 day & 4 & 7,530--21,386 & 0.35 & 0.90 & 0.09 \\
\midrule
ETTh1 & 1 hour & 7 & 17,420 & 0.33 & 0.92 & 0.08 \\
ETTh2 & 1 hour & 7 & 17,420 & 0.61 & 0.92 & 0.07 \\
Weather & 10 min & 21 & 52,704 & 0.78 & 0.82 & 0.11 \\
Electricity & 1 hour & 321 & 26,304 & 0.41 & 0.58 & 0.32 \\
\bottomrule
\end{tabular}
\vspace{1mm}
\parbox{0.98\textwidth}{\scriptsize \textit{Notes.} $M$: numeric variables excluding date; $T$: full-series observations. BRGM $M$/$T$ are ranges and reported STL scores are medians. STL is fitted on the last at most 6,000 observations of the first 50\% of each target series; ADF uses at most 5,000 training observations. BRGM uses an annual STL cycle (365 days); the public series use daily cycles determined from their recorded timestamps. Residual ratio is $\mathrm{Var}(R)/\mathrm{Var}(y)$, not directly observed physical measurement noise. Electricity scores concern target column 319 in the anonymized benchmark CSV, while $M=321$ describes all client series. The original meter identifier for this column has not been verified.}
\end{table*}
